\slajdid{02-s049}
\begin{frame}[label=02-s049]
\gusttekst
Za realizaciju samo sa NILI logičkim kolima po pravilu je najlakše krenuti od proizvoda zbirova
\[
\begin{aligned}
F&=(D+\bar B)(\bar D+\bar C)(\bar D+\bar A)\\
 &=\overline{\overline{(D+\bar B)(\bar D+\bar C)(\bar D+\bar A)}}
 =\overline{\overline{(D+\bar B)}+\overline{(\bar D+\bar C)}+\overline{(\bar D+\bar A)}}
\end{aligned}
\]
pa nam za realizaciju treba pa nam za realizaciju treba čip sa najmanje 4 invertora, 1 čip sa najmanje 3 dvoulazna NILI kola i jedan čip sa najmanje jednim troulaznim NILI kolom. Tri čipa ali različita. I za ovu realizaciju samo sa NILI kolima ne postoji dokaz da je minimalna sa stanovišta klasične minimizacije. Možemo se samo nadati da je tako. Kako se u jednom čipu nalazi 4 dvoulazna NILI kola, možemo ih iskoristiti kao invertore za potrebne 4 inverzije. I isto tako jedno troulazno NILI kolo zameniti sa dva dvoulazna NILI kola i inverzijom
\(\overline{X+Y+Z}=\overline{(X+Y)+Z}=\overline{\bigl(\overline{\overline{X+Y}}\bigr)+Z}\),
ukupan broji dvoulaznih NILI kola koja su nam potrebna je 10. Ukupan broj čipova 3. Ostaje na primer i da se proveri: pod uslovom da posedujemo invertovane vrednosti primenjivih i da ovu kombinacionu mrežu realizujemo kao jednostepeno složeno CMOS kolo koja realizacija bi bila minimalnija sa stanovišta površine koju zauzima. \underline{Probajte}.
\end{frame}
